\subsection*{الف}
برای هموردایی نسبت به انتقال کافی است نشان دهیم جابجایی از کانولوشن خارج می‌شود:
\begin{flalign*}
(T_{a,b}(K*I))(i,j) &= (K*I)(i-a,j-b) \\
&= \sum_{m,n} K(m,n)\, I(i-a-m,\, j-b-n) \\
&= \sum_{m,n} K(m,n)\, (T_{a,b}I)(i-m,\, j-n) \\
&= (K*T_{a,b}I)(i,j).
\end{flalign*}
پس کانولوشن نسبت به انتقال همورداست.

\subsection*{ب}
هر عملگر خطی هموردا نسبت به انتقال، کانولوشن است. کافی است پاسخ عملگر به دلتا را هسته بگیریم:
\begin{flalign*}
K(m,n) = \mathcal{L}(\delta)(m,n), \qquad \delta(i,j)=
\begin{cases}
1 & i=0,\, j=0\\
0 & \text{otherwise}
\end{cases}
\end{flalign*}
با خطی بودن و هموردایی:
\begin{flalign*}
I(i,j) = \sum_{m,n} I(m,n)\, T_{m,n}\delta(i,j) 
\Rightarrow \mathcal{L}(I) = \sum_{m,n} I(m,n)\, T_{m,n}\mathcal{L}(\delta)
= K * I.
\end{flalign*}
نتیجه: طراحی شبکه‌های کانولوشنی دقیقا با اعمال قیود هموردایی (اشتراک وزن‌ها) انجام می‌شود و فضای پارامترها کوچک‌تر و تعمیم بهتر می‌شود.

\subsection*{ج}
اگر تصویر به بردار $x\in \mathbb{R}^{HW}$ تبدیل شود، یک عملگر خطی هموردا نسبت به انتقال باید به صورت ماتریس \lr{Toeplitz} (در دوبعد، \lr{block Toeplitz with Toeplitz blocks}) باشد، چون هر شیفت خروجی باید با همان شیفت در ورودی متناظر شود و ضرایب وابسته به اختلاف اندیس‌ها هستند.  
لایه خطی کاملا متصل \lr{MLP} چنین ساختار ثابتی ندارد و وزن‌ها آزادند، پس به طور کلی هموردایی نسبت به انتقال را تضمین نمی‌کند.

\subsection*{د}
برای تشخیص یک الگوی $k\times k$ در هر مکان:
\begin{flalign*}
\text{کانولوشنی:} &\quad k^2 \text{ وزن } + 1 \text{ بایاس} \;\; \Rightarrow\;\; k^2+1.\\
\text{کاملا متصل:} &\quad \text{برای هر مکان یک الگو لازم است} \\
&\quad N_p=(H-k+1)(W-k+1) \text{ مکان} \\
&\quad \Rightarrow\; (H-k+1)(W-k+1)\,k^2 + (H-k+1)(W-k+1).
\end{flalign*}
پس لایه کانولوشنی به طور چشمگیری پارامترهای کمتری نیاز دارد.